% p.48 — الأفعال المعتلة (تابع): الأعمدة (67)–(72)، باب فَاعَلَ ثم باب تَفَعَّلَ
% Layout as p.41: two blocks per column (أصل above the rule, forms after إعلال below).
% In column (70) the lithograph writes the initial hamza+alif as an alif with a small
% vertical stroke above it (اٰ); it is transcribed here with آ.
\mspage{48}
\begin{sarf}{|c|c|c|c|c|c|}
\thead{(67) مضاعف} & \thead{(68) ناقص واوي} & \thead{(69) ناقص يائي} & \thead{(70) لفيف ومهموز} & \thead{(71) ناقص واوي} & \thead{(72) ناقص يائي}\\\hline
فَاعَلَ & فَاعَلَ & فَاعَلَ & فَاعَلَ & تَفَعَّلَ & تَفَعَّلَ\\\hline
حَابَبَ & رَاضَوَ & لَاقَيَ & \unclear{آوَيَ} & تَخَلَّوَ & تَوَفَّيَ\\
يُحَابِبُ & يُرَاضِوُ & يُلَاقِيُ & \unclear{يُآوِيُ} & يَتَخَلَّوُ & يَتَوَفَّيُ\\
مُحَابَبَةً & مُرَاضَوَةً & مُلَاقَيَةً & \unclear{مُآوَيَةً} & تَخَلُّوًا & تَوَفُّيًا\\
مُحَابِبٌ & مُرَاضِوٌ & مُلَاقِيٌ & \unclear{مُآوِيٌ} & مُتَخَلِّوٌ & مُتَوَفِّيٌ\\
حُوْبِبَ & رُوْضِوَ & لُوْقِيَ & اُوْوِيَ & تُخُلِّوَ & تُوُفِّيَ\\
يُحَابَبُ & يُرَاضَوُ & يُلَاقَيُ & \unclear{يُآوَيُ} & يُتَخَلَّوُ & يُتَوَفَّيُ\\
مُحَابَبَةً & مُرَاضَوَةً & مُلَاقَيَةً & \unclear{مُآوَيَةً} & تَخَلُّوًا & تَوَفُّيًا\\
مُحَابَبٌ & مُرَاضَوٌ & مُلَاقَيٌ & \unclear{مُآوَيٌ} & مُتَخَلَّوٌ & مُتَوَفَّيٌ\\
حَابِبْ & رَاضِوْ & لَاقِيْ & \unclear{آوِيْ} & تَخَلَّوْ & تَوَفَّيْ\\
لَا تُحَابِبْ & لَا تُرَاضِوْ & لَا تُلَاقِيْ & \unclear{لَا تُآوِيْ} & لَا تَتَخَلَّوْ & لَا تَتَوَفَّيْ\\\hline
حَابَّ & رَاضٰى & لَاقٰى & \unclear{آوٰى} & تَخَلّٰى & تَوَفّٰى\\
يُحَابُّ & يُرَاضِيْ & يُلَاقِيْ & \unclear{يُآوِيْ} & يَتَخَلّٰى & يَتَوَفّٰى\\
مُحَابَّةً & مُرَاضَاةً & مُلَاقَاةً & \unclear{مُآوَاةً} & تَخَلِّيًا & تَوَفِّيًا\\
مُحَابٌّ & مُرَاضٍ & مُلَاقٍ & \unclear{مُآوٍ} & مُتَخَلٍّ & مُتَوَفٍّ\\
حُوْبَّ & رُوْضِيَ & لُوْقِيَ & اُوْوِيَ & تُخُلِّيَ & تُوُفِّيَ\\
يُحَابُّ & يُرَاضٰى & يُلَاقٰى & \unclear{يُآوٰى} & يُتَخَلّٰى & يُتَوَفّٰى\\
مُحَابَّةً & مُرَاضَاةً & مُلَاقَاةً & \unclear{مُآوَاةً} & تَخَلِّيًا & تَوَفِّيًا\\
مُحَابٌّ & مُرَاضًى & مُلَاقًى & \unclear{مُآوًى} & مُتَخَلًّى & مُتَوَفًّى\\
حَابَّ & رَاضِ & لَاقِ & \unclear{آوِ} & تَخَلَّ & تَوَفَّ\\
لَا تُحَابَّ & لَا تُرَاضِ & لَا تُلَاقِ & \unclear{لَا تُآوِ} & لَا تَتَخَلَّ & لَا تَتَوَفَّ\\
\end{sarf}
